\documentclass[11pt,a4paper]{article}

% -----------------------
% Layout & typography
% -----------------------
\usepackage[a4paper,margin=1.8cm]{geometry}
\usepackage[T1]{fontenc}
\usepackage[utf8]{inputenc}
\usepackage{lmodern}

% -----------------------
% Recipe environments
% -----------------------
\newenvironment{recipe}[3]{%
  \section*{#1}%
  \textbf{Servings:} #2 \quad \textbf{Appliances:} #3%
  \par\vspace{0.4em}%
}{%
  \vspace{0.8em}%
}

\newenvironment{ingredients}{%
  \subsection*{Ingredients}%
  \begin{itemize}%
}{%
  \end{itemize}%
}

\newenvironment{method}[1][Method]{%
  \subsection*{#1}%
  \begin{enumerate}%
}{%
  \end{enumerate}%
}
 

% -----------------------
% Maths, tables, lists
% -----------------------
\usepackage{amsmath,amssymb}
\usepackage{nicefrac}
\usepackage{booktabs,array,enumitem}
\setlist[itemize]{leftmargin=1.5em}
\setlist[enumerate]{leftmargin=1.8em}

% -----------------------
% Section formatting
% -----------------------
\usepackage{titlesec}
\titleformat{\section}{\large\bfseries}{}{0pt}{}
\titleformat{\subsection}{\normalsize\bfseries}{}{0pt}{}
\titleformat{\paragraph}{\itshape}{}{0pt}{}

% -----------------------
% Spacing & links
% -----------------------
\setlength{\parindent}{0pt}
\setlength{\parskip}{0.5em}
\usepackage[hidelinks]{hyperref}

% -----------------------
% Nutrition table helper
% -----------------------
\newcommand{\nutritiontable}[4]{%
\begin{center}
\begin{tabular}{l r}
\toprule
Energy & #1 kcal \\
Protein & #2 g \\
Fat & #3 g \\
Saturated Fat & #4 g \\
\bottomrule
\end{tabular}
\end{center}
}

% -----------------------
% Document metadata
% -----------------------
\title{Recipe Collection}
\author{}
\date{}

\begin{document}


\begin{recipe}{Chicken, Lentil \& Coconut Curry with Rice}{8}{Instant Pot, Rice Cooker}

\begin{ingredients}
  \item 1.1--1.2\,kg boneless, skinless chicken thighs
  \item 2 cans (400\,g) lentils, drained and rinsed
  \item 2 large onions, diced
  \item 4 garlic cloves, minced
  \item 4 tbsp red curry paste
  \item 1 tsp each ground cumin and ground coriander

  \item 2 cans (400\,ml) light coconut milk
  \item 1 cup chicken stock
  \item 600\,g frozen mixed vegetables
  \item 400\,g frozen broccoli \& cauliflower
  \item 1 tbsp oil
  \item 1 tsp salt
  \item \nicefrac{1}{2}--1 lime, juiced
  \item 520\,g basmati rice (uncooked)
\end{ingredients}

\begin{method}[Method (Instant Pot)]
  \item Set Instant Pot to \textbf{Sauté}. Add oil and cook onion for \textbf{5 minutes} until softened.
  \item Add garlic, red curry paste, cumin, and coriander. Stir for \textbf{1 minute} until fragrant.
  \item Add chicken, coconut milk, stock, and salt. Stir gently to combine.
  \item Pressure cook on \textbf{High} for \textbf{10 minutes}. Allow a \textbf{10-minute natural release}.
  \item Return to \textbf{Sauté}. Stir in lentils and frozen vegetables.
  \item Cook for \textbf{5--7 minutes} until vegetables are tender.
  \item Remove from heat and stir through lime juice. Adjust seasoning.
\end{method}

\begin{method}[Rice (Rice Cooker)]
  \item Rinse rice until water runs mostly clear.
  \item Cook rice with water according to rice cooker instructions.
  \item Fluff and portion (\textbf{~200\,g cooked rice per serve}).
\end{method}

\subsection*{Estimated Nutrition (Per Serve)}
\nutritiontable{680}{40}{17.4}{8.9}

\end{recipe}





\section*{Beef \& Mixed Bean Chilli (Budget-Friendly)}

\textbf{Servings:} 6 \\
\textbf{Appliance:} Instant Pot (5.7 L) \\
\textbf{Prep time:} 10 minutes \\
\textbf{Cook time:} 15 minutes pressure + 10 minutes natural release

\vspace{0.5em}

\subsection*{Ingredients}
\begin{itemize}
    \item Lean beef mince — 900 g
    \item Onion — 1 large, diced
    \item Garlic — 3 cloves, minced
    \item Canned tomatoes — 2 $\times$ 400 g
    \item Kidney beans — 2 $\times$ 400 g, drained and rinsed
    \item Black beans — 1 $\times$ 400 g, drained and rinsed
    \item Frozen mixed vegetables — 600 g
    \item Chilli powder — 2 tbsp
    \item Cumin — 1 tbsp
    \item Paprika — 1 tbsp
    \item Salt — 1$\tfrac{1}{4}$ tsp
    \item Stock or water — 1$\tfrac{1}{2}$ cups
    \item Oil — for sautéing
\end{itemize}

\vspace{0.5em}

\subsection*{Method}
\begin{enumerate}
    \item Set the Instant Pot to \textbf{Sauté}. Add a small amount of oil and cook the onion for 2--3 minutes until softened.
    \item Add the beef mince and cook until well browned, breaking it up with a spoon.
    \item Stir in the garlic, chilli powder, cumin, and paprika. Cook for 30 seconds until fragrant.
    \item Add the canned tomatoes, kidney beans, black beans, frozen mixed vegetables, salt, and stock or water. Stir well, scraping the bottom to prevent sticking.
    \item Secure the lid and cook on \textbf{High Pressure} for 15 minutes.
    \item Allow a \textbf{Natural Pressure Release} for 10 minutes, then vent remaining pressure.
    \item Taste and adjust seasoning. If a thicker chilli is desired, use \textbf{Sauté} mode for a further 5 minutes.
\end{enumerate}

\vspace{0.5em}

\subsection*{Serving Suggestions}
Serve with rice, baked potatoes, or tortillas. Also works well on its own for a lower-carb option. Suitable for meal prep and freezing.

\vspace{0.5em}

\subsection*{Notes}
\begin{itemize}
    \item Frozen mixed vegetables add bulk, fibre, and nutrients at low cost with no additional prep.
    \item Using mixed beans stretches the mince while maintaining protein and satiety.
    \item This recipe freezes well for up to 3 months.
    \item Optional additions include frozen peas, frozen broccoli or cauliflower, or a handful of frozen spinach stirred through when reheating.
\end{itemize}

\section*{Budget Pulled BBQ Chicken with Lentils (8 Servings)}

\textbf{Servings:} 8 \\
\textbf{Appliance:} Instant Pot (5.7 L) \\
\textbf{Prep time:} 10 minutes \\
\textbf{Cook time:} 12 minutes pressure + 10 minutes natural release

\vspace{0.5em}

\subsection*{Why This Works}
\begin{itemize}
    \item High protein at a lower cost per serving
    \item Lentils bulk out the dish without changing flavour
    \item Frozen vegetables add volume and fibre at minimal cost
    \item BBQ sauce masks lentils and vegetables completely
    \item Remains moist and reheats well
\end{itemize}

\vspace{0.5em}

\subsection*{Ingredients}
\begin{itemize}
    \item Chicken breast — 1.1 kg
    \item Red lentils (dry) — 1 cup
    \item Frozen mixed vegetables — 300--400 g
    \item BBQ sauce — 1$\tfrac{1}{4}$ cups (any affordable brand)
    \item Onion — 1 large, thinly sliced
    \item Apple cider vinegar — 2 tbsp
    \item Smoked paprika — 2 tsp
    \item Salt — 1 tsp
    \item Black pepper — $\tfrac{1}{2}$ tsp
    \item Water — 1$\tfrac{1}{2}$ cups
\end{itemize}

\vspace{0.5em}

\subsection*{Prep Alongside}

\paragraph{Rice}
\begin{itemize}
    \item Basmati rice (dry) — 3 cups
    \item Water — 6 cups
    \item Salt — 1 tsp
\end{itemize}

\paragraph{Slaw}
\begin{itemize}
    \item Green cabbage — $\tfrac{1}{2}$ medium head, shredded
    \item Carrots — 2, grated
\end{itemize}

\paragraph{Light Dressing}
\begin{itemize}
    \item Mayo or yoghurt — 3 tbsp
    \item Lemon juice or vinegar — 1 tbsp
    \item Salt and pepper — to taste
\end{itemize}

\vspace{0.5em}

\subsection*{Method}
\begin{enumerate}
    \item Add the sliced onion to the bottom of the Instant Pot.
    \item Sprinkle the red lentils evenly over the onion.
    \item Place the chicken breasts on top, whole.
    \item In a bowl, mix the BBQ sauce, apple cider vinegar, smoked paprika, salt, black pepper, and water.
    \item Pour the sauce mixture evenly over the chicken.
    \item Secure the lid and cook on \textbf{High Pressure} for 12 minutes.
    \item Allow a \textbf{Natural Pressure Release} for 10 minutes, then vent remaining pressure.
    \item Remove the chicken and shred using two forks.
    \item Return the shredded chicken to the pot and stir well to combine with the lentils and sauce.
    \item Stir in the frozen vegetables directly from frozen.
    \item Set the Instant Pot to \textbf{Sauté} and cook for 4--6 minutes, stirring occasionally, until the vegetables are heated through and the sauce thickens slightly.
\end{enumerate}

\vspace{0.5em}

\subsection*{Storage}
\begin{itemize}
    \item 6 servings: refrigerate for up to 4 days
    \item 2 servings: freeze and label \emph{``Week 2 BBQ Chicken -- Fri Dinner''}
\end{itemize}

\vspace{0.5em}

\subsection*{Serving Ideas}
\begin{itemize}
    \item \textbf{Lunch bowl:} Pulled BBQ chicken, $\tfrac{3}{4}$ cup rice, generous slaw
    \item \textbf{Light dinner:} Pulled BBQ chicken, $\tfrac{1}{2}$ cup rice or wrap, extra slaw or vegetables
\end{itemize}

\vspace{0.5em}

\subsection*{Notes}
\begin{itemize}
    \item Frozen vegetables increase meal volume with minimal added cost or calories.
    \item Add vegetables after pressure cooking to preserve texture.
    \item Lentils fully soften and blend into the sauce, increasing protein and volume.
    \item Chicken thigh may be substituted for breast for additional savings and moisture.
    \item Recipe freezes well with no texture loss.
\end{itemize}

\section*{Budget Teriyaki Chicken \& Veg (6 Servings)}

\textbf{Servings:} 6 \\
\textbf{Appliance:} Instant Pot (5.7 L) \\
\textbf{Prep time:} 10 minutes \\
\textbf{Cook time:} 10 minutes pressure + sauté finish

\vspace{0.5em}

\subsection*{Why This Works}
\begin{itemize}
    \item Balanced sweet--savory flavour
    \item Chicken remains the star
    \item Edamame adds affordable protein without dominating
    \item Frozen vegetables keep cost and waste low
\end{itemize}

\vspace{0.5em}

\subsection*{Ingredients}
\begin{itemize}
    \item Chicken thighs (boneless) — 800 g
    \item Frozen shelled edamame — 225 g (half of a 454 g bag)
    \item Teriyaki sauce — $\tfrac{3}{4}$ cup
    \item Soy sauce — 2 tbsp
    \item Honey or brown sugar — 1 tbsp
    \item Garlic — 3 cloves, minced
    \item Ginger (optional) — 1 tsp, grated
    \item Frozen vegetables — 700--800 g
    \item Water — $\tfrac{1}{4}$ cup
\end{itemize}

\vspace{0.5em}

\subsection*{Frozen Vegetable Options}
\begin{itemize}
    \item Oriental or Asian stir-fry mix
    \item Broccoli \& cauliflower mix
    \item Mixed vegetables
    \item Green beans or baby beans
    \item Capsicum mix
\end{itemize}

\vspace{0.5em}

\subsection*{Prep Alongside}

\paragraph{Rice}
\begin{itemize}
    \item Jasmine or basmati rice (dry) — 2 cups
    \item Water — 4 cups
    \item Salt — $\tfrac{1}{2}$ tsp
\end{itemize}

\paragraph{Optional Garnishes}
\begin{itemize}
    \item Sesame seeds
    \item Spring onions
\end{itemize}

\vspace{0.5em}

\subsection*{Method}
\begin{enumerate}
    \item Place the chicken thighs in the bottom of the Instant Pot.
    \item In a bowl, mix the teriyaki sauce, soy sauce, honey or sugar, garlic, ginger, and water.
    \item Pour the sauce evenly over the chicken.
    \item Secure the lid and cook on \textbf{High Pressure} for 10 minutes.
    \item Perform a \textbf{Quick Release} of pressure.
    \item Add the frozen vegetables and frozen edamame directly to the pot.
    \item Set to \textbf{Sauté} and cook for 5--7 minutes, stirring gently, until vegetables are tender and the sauce thickens slightly.
    \item Taste and adjust seasoning if needed.
\end{enumerate}

\vspace{0.5em}

\subsection*{Storage}
\begin{itemize}
    \item Refrigerate for up to 4 days
    \item Freezes well for up to 3 months
\end{itemize}

\vspace{0.5em}

\subsection*{Serving Ideas}
\begin{itemize}
    \item Serve over rice for a balanced bowl
    \item Use leftovers in wraps or lettuce cups
    \item Add extra vegetables to stretch further
\end{itemize}

\vspace{0.5em}

\subsection*{Notes}
\begin{itemize}
    \item Half a bag of edamame provides protein without overwhelming texture.
    \item Remaining edamame can be used in fried rice or stir-fries.
    \item Chicken thigh is preferred for cost and moisture.
\end{itemize}

\section*{Budget Beef \& Vegetable Stew (8 Servings)}

\textbf{Servings:} 8 \\
\textbf{Appliance:} Instant Pot (5.7 L) \\
\textbf{Prep time:} 10 minutes \\
\textbf{Cook time:} 35 minutes pressure + 15 minutes natural release

\vspace{0.5em}

\subsection*{Why This Works}
\begin{itemize}
    \item Extremely filling: protein, fibre, and volume
    \item Built-in carbohydrates (no sides required)
    \item Lentils stretch the beef without changing the flavour
    \item Frozen vegetables reduce cost, prep time, and waste
    \item Improves after resting overnight
    \item Ideal cold-weather, low-effort meal
\end{itemize}

\vspace{0.5em}

\subsection*{Ingredients}
\begin{itemize}
    \item Beef chuck — 1.2 kg, cut into large cubes
    \item Red lentils (dry) — 1 cup
    \item Frozen mixed vegetables — 1 kg (peas, corn, carrot, beans)
    \item Potatoes — 800 g, diced into 2--3 cm chunks
    \item Onion — 2 large, diced
    \item Garlic — 4 cloves, minced
    \item Beef stock — 3 cups (750 ml)
    \item Tomato paste — 3 tbsp
    \item Worcestershire sauce — 2 tbsp
    \item Salt — 2 tsp
    \item Black pepper — 1 tsp
    \item Dried thyme or mixed herbs — 2 tsp
    \item Olive oil — 1 tbsp
\end{itemize}

\vspace{0.5em}

\subsection*{Method}
\begin{enumerate}
    \item Set the Instant Pot to \textbf{Sauté}. Add the olive oil.
    \item Cook the onion for approximately 5 minutes until softened.
    \item Add the garlic and cook for 30 seconds until fragrant.
    \item Add the beef cubes and stir briefly to coat in the onion mixture.
    \item Add the red lentils and diced potatoes.
    \item Stir in the tomato paste, Worcestershire sauce, herbs, salt, and pepper.
    \item Pour in the beef stock and mix well, scraping the bottom of the pot.
    \item Cancel \textbf{Sauté}. Secure the lid with the valve sealed.
    \item Cook on \textbf{High Pressure} for 35 minutes.
    \item Allow a \textbf{Natural Pressure Release} for 15 minutes, then vent remaining pressure.
    \item Open the lid and stir in the frozen mixed vegetables.
    \item Set to \textbf{Sauté} and simmer uncovered for 5--10 minutes until vegetables are heated through and the stew thickens slightly.
    \item Taste and adjust seasoning if required.
\end{enumerate}

\vspace{0.5em}

\subsection*{Storage}
\begin{itemize}
    \item Refrigerate for up to 4 days
    \item Freeze portions for up to 3 months
    \item Flavour improves after 24 hours
\end{itemize}

\vspace{0.5em}

\subsection*{Serving Notes}
\begin{itemize}
    \item Serve as-is; potatoes and lentils provide sufficient carbohydrates
    \item Excellent with bread if extra calories are needed
    \item Ideal for reheating on low heat
\end{itemize}

\vspace{0.5em}

\subsection*{Notes}
\begin{itemize}
    \item Lentils dissolve into the broth, thickening the stew naturally.
    \item Beef shin or gravy beef can be substituted for chuck if cheaper.
    \item Frozen broccoli or cauliflower can replace up to 300 g of the mixed vegetables if preferred.
\end{itemize}

\section*{Budget Lentil Coconut Curry (6 Servings)}

\textbf{Servings:} 6 \\
\textbf{Appliance:} Instant Pot (5.7 L) \\
\textbf{Prep time:} 10 minutes \\
\textbf{Cook time:} 12 minutes pressure + 10 minutes natural release

\vspace{0.5em}

\subsection*{Why This Works}
\begin{itemize}
    \item Very low cost per serving
    \item High protein and fibre from lentils
    \item Filling but easy to digest
    \item Excellent meat-free contrast to heavier stews
    \item Ideal for lunches and reheats well
    \item Frozen vegetables keep cost low and eliminate waste
\end{itemize}

\vspace{0.5em}

\subsection*{Ingredients}
\begin{itemize}
    \item Dried red lentils — 2$\tfrac{1}{2}$ cups, rinsed
    \item Onion — 1 large, diced
    \item Garlic — 3 cloves, minced
    \item Curry powder — 2 tbsp
    \item Ground cumin — 1 tbsp
    \item Canned tomatoes — 2 $\times$ 400 g
    \item Light coconut milk — 1 $\times$ 400 ml
    \item Water — 2 cups (500 ml)
    \item Salt — 1$\tfrac{1}{2}$ tsp
    \item Frozen vegetables — 1--2 cups (see options below)
\end{itemize}

\vspace{0.5em}

\subsubsection*{Recommended Frozen Vegetable Options}
\begin{itemize}
    \item Pam's Mixed Vegetables (peas, corn, carrot, beans) — best value and texture
    \item Broccoli \& cauliflower mix — higher fibre, slightly firmer texture
    \item Garden peas — optional for added bulk and mild sweetness
    \item Frozen spinach — optional; use as a nutrient boost rather than bulk
\end{itemize}

\vspace{0.5em}

\subsection*{Method}
\begin{enumerate}
    \item Add the lentils, onion, garlic, curry powder, and ground cumin to the Instant Pot.
    \item Add the canned tomatoes, coconut milk, water, and salt.
    \item Stir well to combine, ensuring lentils are fully submerged.
    \item Secure the lid with the valve sealed and cook on \textbf{High Pressure} for 12 minutes.
    \item Allow a \textbf{Natural Pressure Release} for 10 minutes, then vent remaining pressure.
    \item Open the lid and stir thoroughly; the lentils will thicken rapidly.
    \item Stir in frozen vegetables.
    \item Select \textbf{Sauté} mode and cook for 5--7 minutes until vegetables are heated through.
    \item Adjust thickness with additional water if required.
\end{enumerate}

\vspace{0.5em}

\subsection*{Portioning \& Storage}
\begin{itemize}
    \item All 6 servings: refrigerate
    \item Best consumed within 4 days
    \item Thickens further after chilling; thin with water when reheating
\end{itemize}

\vspace{0.5em}

\subsection*{What to Prep Alongside}

\paragraph{Rice}
\begin{itemize}
    \item Dry rice — 2$\tfrac{1}{2}$ cups
    \item Yield — approximately 7 cups cooked
\end{itemize}

\paragraph{Vegetables}
\begin{itemize}
    \item No additional vegetables required; curry is designed as a one-pot meal
    \item Optional: extra steamed frozen veg for larger portions
\end{itemize}

\vspace{0.5em}

\subsection*{Serving Ideas}
\begin{itemize}
    \item \textbf{Lunch:} Lentil coconut curry, $\tfrac{3}{4}$ cup rice
    \item \textbf{Dinner:} Lentil coconut curry, $\tfrac{1}{2}$ cup rice, extra vegetables if desired
\end{itemize}

\vspace{0.5em}

\subsection*{Notes}
\begin{itemize}
    \item Frozen pumpkin or diced potatoes (2--3 cups) may be added for extra bulk without changing cook time.
    \item Full-fat coconut milk can be used for richer flavour if budget allows.
    \item This curry freezes well; texture will be thicker after thawing.
\end{itemize}

\section*{Budget Chicken \& Vegetable Soup (8 Servings)}

\textbf{Servings:} 8 \\
\textbf{Appliance:} Instant Pot (5.7 L) \\
\textbf{Prep time:} 10 minutes \\
\textbf{Cook time:} 15 minutes pressure + 10 minutes natural release

\vspace{0.5em}

\subsection*{Why This Works}
\begin{itemize}
    \item High protein and very high volume
    \item Frozen vegetables provide cheap bulk with zero waste
    \item Lentils add low-cost protein and natural thickness
    \item Minimal prep and no browning required
    \item Excellent for early-week lunches and dinners
    \item Light but filling
\end{itemize}

\vspace{0.5em}

\subsection*{Ingredients}

\paragraph{Protein}
\begin{itemize}
    \item Chicken thighs or breast — 1.2 kg
\end{itemize}

\paragraph{Vegetables}
\begin{itemize}
    \item Frozen mixed vegetables (peas, corn, carrot, beans) — 1 kg
    \item Optional frozen broccoli \& cauliflower — 300--500 g
    \item Onion — 1 large, diced
    \item Garlic — 3 cloves, minced
\end{itemize}

\paragraph{Bulk \& Protein Extender}
\begin{itemize}
    \item Red lentils (dry) — 1 cup, rinsed
\end{itemize}

\paragraph{Liquid \& Seasoning}
\begin{itemize}
    \item Chicken stock — 3.5 L
    \item Salt — 2 tsp (adjust to taste)
    \item Black pepper — 1 tsp
    \item Dried mixed herbs — 2 tsp
    \item Bay leaves — 2 (optional)
\end{itemize}

\vspace{0.5em}

\subsection*{Method}
\begin{enumerate}
    \item Add all ingredients to the Instant Pot. Place the chicken in whole.
    \item Stir gently to distribute lentils and vegetables.
    \item Secure the lid with the valve sealed.
    \item Cook on \textbf{High Pressure} for 15 minutes.
    \item Allow a \textbf{Natural Pressure Release} for 10 minutes, then vent remaining pressure.
    \item Remove the chicken and shred using two forks.
    \item Return the shredded chicken to the pot and stir well.
    \item Taste and adjust salt and pepper as required.
    \item If the soup is too thick, add additional water or stock.
\end{enumerate}

\vspace{0.5em}

\subsection*{Portioning \& Storage}
\begin{itemize}
    \item Portion into 8 equal containers
    \item Refrigerate 6 servings for the week
    \item Freeze 2 servings for later use (e.g.\ Friday dinners)
\end{itemize}

\vspace{0.5em}

\subsection*{Serving Notes}
\begin{itemize}
    \item Serve as-is for a light meal
    \item Add bread or toast if extra calories are needed
    \item Thickens slightly after chilling; thin when reheating
\end{itemize}

\vspace{0.5em}

\subsection*{Notes}
\begin{itemize}
    \item Frozen vegetables are added from frozen; no thawing required.
    \item Red lentils soften completely and blend into the broth.
    \item Chicken thighs provide more flavour; breast is leaner but still suitable.
    \item Optional frozen spinach can be stirred through when reheating individual portions.
\end{itemize}

\section*{Budget Protein Pasta Sauce (6 Servings)}

\textbf{Servings:} 6 \\
\textbf{Appliance:} Instant Pot (5.7 L) \\
\textbf{Prep time:} 10 minutes \\
\textbf{Cook time:} 15 minutes pressure + 10 minutes natural release

\vspace{0.5em}

\subsection*{Why This Works}
\begin{itemize}
    \item High protein without heavy meals
    \item Flexible: meat-based or plant-based
    \item Frozen vegetables add volume and fibre at low cost
    \item Ideal for structured meal prep and budget control
    \item Reheats cleanly and pairs with fresh pasta
\end{itemize}

\vspace{0.5em}

\subsection*{Ingredients}

\paragraph{Protein (Choose One)}
\begin{itemize}
    \item Lean beef mince — 800 g \emph{(budget meat option)}
    \item Lean beef mince — 600 g + dried red lentils — 1 cup \emph{(budget extender)}
\end{itemize}

\textit{OR plant-based option:}

\begin{itemize}
    \item Dried red lentils — 2$\tfrac{1}{2}$ cups
\end{itemize}

\paragraph{Vegetables (Frozen, Budget-Optimised)}
\begin{itemize}
    \item Frozen mixed vegetables — 600 g
    \item Frozen broccoli \& cauliflower mix — 300 g
    \item Onion — 1 large, diced
    \item Garlic — 3 cloves, minced
\end{itemize}

\paragraph{Tomato Base}
\begin{itemize}
    \item Canned crushed tomatoes — 2 $\times$ 400 g
    \item Tomato paste — 3 tbsp
    \item Water — 1 cup
\end{itemize}

\paragraph{Seasoning}
\begin{itemize}
    \item Italian herbs — 2 tbsp
    \item Salt — 1$\tfrac{1}{2}$ tsp
    \item Black pepper — 1 tsp
\end{itemize}

\vspace{0.5em}

\subsection*{Method}
\begin{enumerate}
    \item Set the Instant Pot to \textbf{Sauté}. Cook the onion and garlic for 2--3 minutes until softened.
    \item If using beef mince, add it now and brown lightly. (Skip this step if using lentils only.)
    \item Add lentils (if using), frozen vegetables, crushed tomatoes, tomato paste, water, herbs, salt, and pepper.
    \item Stir well, scraping the bottom of the pot to prevent sticking.
    \item Secure the lid with the valve sealed and cook on \textbf{High Pressure} for 15 minutes.
    \item Allow a \textbf{Natural Pressure Release} for 10 minutes, then vent remaining pressure.
    \item Stir thoroughly. Adjust thickness as required:
    \begin{itemize}
        \item If too thick: add water
        \item If too thin: use \textbf{Sauté} mode for 5 minutes, lid off
    \end{itemize}
\end{enumerate}

\vspace{0.5em}

\subsection*{Pasta Prep (Cook Fresh)}
\begin{itemize}
    \item Dry pasta — 600 g total for the week
    \item Cook per meal to prevent over-softening
\end{itemize}

\vspace{0.5em}

\subsection*{Portioning \& Storage}
\begin{itemize}
    \item Portion sauce into 6 containers
    \item Refrigerate up to 4 days or freeze individual portions
\end{itemize}

\vspace{0.5em}

\subsection*{How to Eat It}
\begin{itemize}
    \item \textbf{Lunch:} Sauce with 75--90 g dry pasta, plus extra vegetables if desired
    \item \textbf{Dinner:} Sauce with 50--75 g dry pasta, plus a large vegetable portion
\end{itemize}

\vspace{0.5em}

\subsection*{Notes}
\begin{itemize}
    \item Frozen vegetables reduce cost, waste, and prep time while maintaining volume.
    \item Lentils soften completely and blend into the sauce, increasing protein and fibre.
    \item This sauce freezes well without texture loss.
    \item Lower pasta portions at dinner support lighter digestion and improved sleep.
\end{itemize}


\end{document}
